%%
%% This is file `sample-manuscript.tex',
%% generated with the docstrip utility.
%%
%% The original source files were:
%%
%% samples.dtx  (with options: `all,proceedings,bibtex,manuscript')
%% 
%% IMPORTANT NOTICE:
%% 
%% For the copyright see the source file.\textbf{}
%% 
%% Any modified versions of this file must be renamed
%% with new filenames distinct from sample-manuscript.tex.
%% 
%% For distribution of the original source see the terms
%% for copying and modification in the file samples.dtx.
%% 
%% This generated file may be distributed as long as the
%% original source files, as listed above, are part of the
%% same distribution. (The sources need not necessarily be
%% in the same archive or directory.)
%%
%%
%% Commands for TeXCount
%TC:macro \cite [option:text,text]
%TC:macro \citep [option:text,text]
%TC:macro \citet [option:text,text]
%TC:envir table 0 1
%TC:envir table* 0 1
%TC:envir tabular [ignore] word
%TC:envir displaymath 0 word
%TC:envir math 0 word
%TC:envir comment 0 0
%%
%% The first command in your LaTeX source must be the \documentclass
%% command.
%%
%% For submission and review of your manuscript please change the
%% command to \documentclass[manuscript, screen, review]{acmart}.
%%
%% When submitting camera ready or to TAPS, please change the command
%% to \documentclass[sigconf]{acmart} or whichever template is required
%% for your publication.
%%
%%
\documentclass[manuscript,screen,review]{acmart}
%%
%% \BibTeX command to typeset BibTeX logo in the docs
\AtBeginDocument{%
  \providecommand\BibTeX{{%
    Bib\TeX}}}

%% Rights management information.  This information is sent to you
%% when you complete the rights form.  These commands have SAMPLE
%% values in them; it is your responsibility as an author to replace
%% the commands and values with those provided to you when you
%% complete the rights form.
\setcopyright{acmlicensed}
\copyrightyear{2018}
\acmYear{2018}
\acmDOI{XXXXXXX.XXXXXXX}
%% These commands are for a PROCEEDINGS abstract or paper.
\acmConference[Conference acronym 'XX]{Make sure to enter the correct
  conference title from your rights confirmation email}{June 03--05,
  2018}{Woodstock, NY}
%%
%%  Uncomment \acmBooktitle if the title of the proceedings is different
%%  from ``Proceedings of ...''!
%%
%%\acmBooktitle{Woodstock '18: ACM Symposium on Neural Gaze Detection,
%%  June 03--05, 2018, Woodstock, NY}
\acmISBN{978-1-4503-XXXX-X/2018/06}


%%
%% Submission ID.
%% Use this when submitting an article to a sponsored event. You'll
%% receive a unique submission ID from the organizers
%% of the event, and this ID should be used as the parameter to this command.
%%\acmSubmissionID{123-A56-BU3}

%%
%% For managing citations, it is recommended to use bibliography
%% files in BibTeX format.
%%
%% You can then either use BibTeX with the ACM-Reference-Format style,
%% or BibLaTeX with the acmnumeric or acmauthoryear sytles, that include
%% support for advanced citation of software artefact from the
%% biblatex-software package, also separately available on CTAN.
%%
%% Look at the sample-*-biblatex.tex files for templates showcasing
%% the biblatex styles.
%%

%%
%% The majority of ACM publications use numbered citations and
%% references.  The command \citestyle{authoryear} switches to the
%% "author year" style.
%%
%% If you are preparing content for an event
%% sponsored by ACM SIGGRAPH, you must use the "author year" style of
%% citations and references.
%% Uncommenting
%% the next command will enable that style.
%%\citestyle{acmauthoryear}


%%
%% end of the preamble, start of the body of the document source.
\begin{document}

%%
%% The "title" command has an optional parameter,
%% allowing the author to define a "short title" to be used in page headers.
\title{Proyecto IOT}

%%
%% The "author" command and its associated commands are used to define
%% the authors and their affiliations.
%% Of note is the shared affiliation of the first two authors, and the
%% "authornote" and "authornotemark" commands
%% used to denote shared contribution to the research.

\author{Luis Antonio Guti\'errez Guanilo}
% \authornote{Both authors contributed equally to this research.}
\email{luis.gutierrez.g@utec.edu.pe}
% \orcid{1234-5678-9012}
% \author{G.K.M. Tobin}
% \authornotemark[1]
% \email{webmaster@marysville-ohio.com}
\affiliation{%
  \institution{University of Engineering and Technology}
  \city{Barranco}
  \state{Lima}
  \country{Peru}
}

\author{Jeremy Jeffrey Matos Cangalaya}
\affiliation{%
  \institution{University of Engineering and Technology}
  \city{Barranco}
  \country{Peru}
}
\email{jeremy.matos.@utec.edu.pe}

\author{Plinio Matias Avenda\~no Vargas}
\affiliation{%
  \institution{University of Engineering and Technology}
  \city{Barranco}
  \country{Peru}
}

\author{Juan Sebastian Sara Junco}
\affiliation{%
 \institution{University of Engineering and Technology}
  \city{Barranco}
  \state{Lima}
  \country{Peru}
}

%%
%% By default, the full list of authors will be used in the page
%% headers. Often, this list is too long, and will overlap
%% other information printed in the page headers. This command allows
%% the author to define a more concise list
%% of authors' names for this purpose.
\renewcommand{\shortauthors}{Guti\'errez et al.}

%%
%% The abstract is a short summary of the work to be presented in the
%% article.
\begin{abstract}
This report explores the development of wearable devices designed to address the growing concern of sedentary lifestyles and poor posture, particularly in Peru and Latin America. With the increasing reliance on digital technologies and long hours spent in sedentary activities, posture-related health issues such as back pain and cardiovascular diseases have become prevalent. The proposed project focuses on the design and implementation of two wearable devices that monitor posture and vital signs, specifically oxygenation and heart rate, in real-time. These devices, integrated with a mobile application, aim to provide users with immediate feedback to correct posture and track their health metrics. The methodology includes the use of sensors such as the MPU6050 for posture detection and the MAX30102 for pulse oximetry, along with the ESP32 microcontroller to process and transmit data wirelessly. The project aims to contribute to the prevention of health problems associated with sedentary behavior and to promote healthier lifestyles.
\end{abstract}


%%
%% The code below is generated by the tool at http://dl.acm.org/ccs.cfm.
%% Please copy and paste the code instead of the example below.
%%
\begin{CCSXML}
<ccs2012>
   <concept>
       <concept_id>10003120.10003121.10003128.10011755</concept_id>
       <concept_desc>Human-centered computing~Gestural input</concept_desc>
       <concept_significance>500</concept_significance>
       </concept>
 </ccs2012>
\end{CCSXML}

\ccsdesc[500]{Human-centered computing~Gestural input}
%%
%% Keywords. The author(s) should pick words that accurately describe
%% the work being presented. Separate the keywords with commas.
\keywords{Do, Not, Us, This, Code, Put, the, Correct, Terms, for,
  Your, Paper}

\received{20 February 2007}
\received[revised]{12 March 2009}
\received[accepted]{5 June 2009}

%%
%% This command processes the author and affiliation and title
%% information and builds the first part of the formatted document.
\maketitle


\section{Introducción}

\subsection{Problemática}
En los últimos años, la vida sedentaria se ha convertido en una de las principales preocupaciones de salud pública, tanto en Perú como en el resto de Latinoamérica. La rápida urbanización, el crecimiento del uso de tecnologías digitales y las largas horas dedicadas a actividades estáticas como el trabajo frente a pantallas, el uso de teléfonos móviles y el consumo de entretenimiento en dispositivos electrónicos han generado un incremento significativo en la inactividad física. Esta falta de movimiento está directamente vinculada a una serie de problemas de salud, entre ellos, dolores musculares, trastornos en la postura, enfermedades cardiovasculares y metabólicas.

En Perú, un estudio realizado por el Ministerio de Salud en 2019 reveló que más del 50\% de la población adulta lleva una vida sedentaria, lo que significa que no cumplen con las recomendaciones mínimas de actividad física diaria. Esta tendencia se ve reflejada en el aumento de enfermedades relacionadas con el sedentarismo, como los dolores musculares y problemas de espalda. Según la Organización Mundial de la Salud (OMS), la mala postura es una de las principales causas de dolor lumbar, un problema que afecta al 70\% de los peruanos, especialmente entre aquellos que pasan más de 4 horas al día sentados o frente a una pantalla. El uso prolongado de dispositivos electrónicos y las malas posturas adoptadas mientras se realiza trabajo o estudio son factores que contribuyen directamente a la aparición de dolor y molestias en la columna vertebral, lo que puede llevar a problemas más graves si no se corrigen a tiempo.

Además, el sedentarismo no solo afecta la postura, sino también la salud cardiovascular. La falta de actividad física regular incrementa el riesgo de enfermedades del corazón, hipertensión y otras condiciones relacionadas con el sistema circulatorio. De acuerdo con la Organización Panamericana de la Salud (OPS), las enfermedades cardiovasculares son la principal causa de muerte en Latinoamérica, y se estima que un alto porcentaje de la población presenta factores de riesgo, como la hipertensión y el colesterol elevado, exacerbados por la falta de ejercicio físico. En Perú, el 29\% de la población adulta padece de hipertensión, una cifra alarmante que refleja la necesidad urgente de adoptar hábitos de vida más saludables, entre ellos, el aumento de la actividad física.

Ante esta problemática, surge la necesidad de desarrollar soluciones tecnológicas que fomenten un cambio en el estilo de vida y ayuden a las personas a controlar su salud de manera más efectiva. El proyecto que se presenta busca ofrecer herramientas accesibles, como un collar inteligente que evalúa la postura y un dispositivo de medición de oxigenación y pulso, para monitorear en tiempo real la salud de los usuarios. Estos dispositivos, conectados a una aplicación móvil, permitirían detectar problemas de postura y salud cardiovascular de manera temprana, fomentando hábitos más saludables y la prevención de enfermedades relacionadas con el sedentarismo.

\section{Marco Te\'orico}

\subsection{Internet de las Cosas (IoT)}
En términos generales, el Internet de las Cosas se entiende como la interconexión de objetos físicos provistos de sensores, capacidad de procesamiento y conectividad de red, los cuales pueden intercambiar datos sin intervención humana directa. Esta noción suele representarse mediante una arquitectura de tres capas: \textit{percepción}, \textit{red} y \textit{aplicación}. Cada capa cumple un rol específico: la primera captura señales del entorno; la segunda las transporta, y la tercera las procesa y presenta al usuario. \cite{kitrum2023}

\subsection{Aplicaciones de los wearables en salud}
Los dispositivos vestibles amplían el paradigma IoT hacia el cuerpo humano. En salud preventiva, permiten monitorear variables fisiológicas y biomecánicas de forma continua, facilitando la detección temprana de riesgo postural, fatiga o hipoxemia. Estudios recientes demuestran que los sensores inerciales integrados en prendas o cinturones alcanzan precisiones superiores al 90\% en la clasificación de posturas de la columna, mientras que los pulsioxímetros ópticos añaden contexto sobre la condición cardiorrespiratoria del usuario.

\subsection{Microcontrolador ESP32 DEVKIT V1}
El núcleo del dispositivo es el ESP32, un SoC de doble procesador Xtensa LX6 a 240 MHz que incorpora 520 kB de SRAM, Wi-Fi 802.11 b/g/n (hasta 150 Mb/s) y Bluetooth clásico/BLE. Entre sus ventajas destaca el modo \emph{deep-sleep} —menor a 10 µA—, esencial para prolongar la autonomía en aplicaciones portátiles. Adicionalmente, la plataforma dispone de aceleradores criptográficos (AES, SHA-2, RSA) y de funciones de arranque seguro y cifrado de memoria flash, lo que fortalece la protección de firmware y de datos biométricos.

\subsection{Sensores de adquisición de datos}

\subsubsection{Sensor inercial MPU6050}
El MPU6050 unifica en un solo encapsulado un acelerómetro y un giroscopio de tres ejes con convertidores A/D de 16 bits. Gracias a su bus I²C y a la posibilidad de fusionar datos mediante filtros complementarios o de Kalman, el módulo permite estimar la orientación (pitch, roll, yaw) con la precisión necesaria para detectar desviaciones posturales en tiempo real.

\subsubsection{Fotopletismografía y Pulsioxímetro MAX30102}
La fotopletismografía (PPG) aprovecha la absorción diferencial de luz en el tejido para inferir cambios en el volumen sanguíneo. El MAX30102 integra LEDs rojo (660 nm) e infrarrojo (880 nm) junto con un fotodiodo y electrónica de rechazo de luz ambiente. Mediante la relación entre las señales AC y DC en ambos canales se calcula la saturación de oxígeno (SpO${_2}$), mientras que la periodicidad de la señal pulsátil revela la frecuencia cardíaca.

\subsection{Gestión energética y protección}
A fin de garantizar autonomía y seguridad, el sistema emplea una batería plana de 3,7 V / 450 mAh junto con el cargador TP4056, el cual implementa carga CC/CV a 1 A y un circuito DW01A que protege frente a sobrecarga, sobredescarga y cortocircuitos. Este conjunto extiende la vida útil de la batería y mitiga riesgos térmicos o eléctricos propios de plataformas portátiles. 

\section{Estado del Arte}

\subsection{A Skin-Integrated Device for Neck Posture Monitoring and Correction}
\textbf{Referencia:} Luo, Hu, et al. (2023) \cite{luo2023skin}\\

\textbf{¿Qué propusieron?} \\
Este artículo presenta el desarrollo de un dispositivo de corrección de la postura del cuello basado en piel electrónica (e-skin). El objetivo es monitorear la postura del cuello en tiempo real y corregirla mediante retroalimentación táctil, con el fin de prevenir trastornos musculoesqueléticos como la espondilosis cervical, común entre las personas que pasan mucho tiempo en posturas estáticas.

\textbf{¿Cómo lo hicieron?} \\
El dispositivo está compuesto por sensores de flexión y actuadores de vibración, integrados en una estructura flexible que se adhiere a la piel del cuello. Utilizaron tecnologías de electrónica flexible para crear un sistema conformable que detecta la inclinación del cuello y proporciona retroalimentación háptica. Este dispositivo fue probado en varios escenarios de uso cotidiano, evaluando su efectividad en la corrección postural durante actividades sedentarias y dinámicas.

\textbf{Resultados:} \\
Los resultados mostraron que el dispositivo pudo monitorear con precisión la postura del cuello y proporcionar retroalimentación efectiva para corregirla. Los usuarios experimentaron una mejora significativa en la alineación postural del cuello, con una reducción notable en el dolor cervical reportado después de usar el dispositivo de manera continua. El sistema demostró ser eficiente en el contexto de la vida cotidiana, con una adherencia robusta y confiable al cuello.

\subsection{PosePilot: An Edge-AI Solution for Posture Correction in Physical Exercises}
\textbf{Referencia:} Gadhvi, Rushiraj, et al. (2025) \cite{gadhvi2025posepilot}\\

\textbf{¿Qué propusieron?} \\
PosePilot es una solución basada en inteligencia artificial (IA) en el borde, diseñada para corregir la postura durante ejercicios físicos, con énfasis en la corrección postural durante la práctica de yoga. El sistema utiliza modelos de aprendizaje profundo para ofrecer retroalimentación en tiempo real sobre la postura del usuario, ayudando a mejorar su rendimiento y prevenir lesiones.

\textbf{¿Cómo lo hicieron?} \\
Utilizaron redes neuronales profundas, específicamente LSTM y BiLSTM con atención multi-cabeza, para procesar los datos de movimiento recogidos de sensores de postura. Los modelos entrenados fueron implementados en dispositivos de borde, permitiendo la corrección postural sin necesidad de conexión a la nube. La solución fue probada en ejercicios de yoga, con un enfoque en la alineación de la postura y la sincronización de los movimientos.

\textbf{Resultados:} \\
El sistema demostró ser altamente preciso en la corrección postural en tiempo real, logrando una mejora en la postura de los usuarios en aproximadamente un 20\% tras la implementación de las recomendaciones del sistema. Los participantes también reportaron una experiencia más satisfactoria y segura al practicar ejercicios físicos, destacando la efectividad de la retroalimentación instantánea proporcionada por la inteligencia artificial.

\subsection{A Smart Chair for Health Monitoring in Daily Life}
\textbf{Referencia:} Nguyen Thi Minh Huong, et al. (2024) \cite{nguyen2024smart}\\

\textbf{¿Qué propusieron?} \\
Este estudio presenta una silla inteligente equipada con sensores de presión y un módulo PPG para monitorear la postura y los signos vitales de los usuarios en tiempo real, especialmente para personas con estilos de vida sedentarios. El sistema no solo clasifica las posturas del cuerpo, sino que también monitorea la frecuencia cardíaca, brindando retroalimentación instantánea sobre la salud del usuario.

\textbf{¿Cómo lo hicieron?} \\
La silla incorpora diez sensores de presión distribuidos en el cojín, junto con un módulo PPG colocado en el brazo del usuario. Utilizaron técnicas de aprendizaje automático para procesar los datos de estos sensores y clasificar las posturas en tiempo real. Los datos obtenidos se mostraban en una aplicación móvil, permitiendo al usuario monitorear su postura y salud de manera continua.

\textbf{Resultados:} \\
El sistema alcanzó una precisión de clasificación de posturas de un 99\% y fue efectivo para reducir el tiempo sedentario. Además, el sistema mostró ser útil en la prevención de problemas musculoesqueléticos y cardiovasculares al proporcionar retroalimentación inmediata sobre la postura y los signos vitales. Los usuarios reportaron mejoras en su bienestar general y una mayor conciencia de sus hábitos posturales.

\subsection{A Rapid Review on the Effectiveness and Use of Wearable Biofeedback Motion Capture Systems in Ergonomics to Mitigate Adverse Postures and Movements of the Upper Body}
\textbf{Referencia:} Lind, Carl M. (2024) \cite{lind2024rapid}\\

\textbf{¿Qué propusieron?} \\
Este artículo revisa la efectividad de los sistemas de captura de movimiento con retroalimentación biofísica utilizados para mitigar posturas y movimientos adversos en la parte superior del cuerpo en contextos laborales. Los sistemas de retroalimentación biofísica son una herramienta clave en la ergonomía para prevenir lesiones musculoesqueléticas asociadas con el trabajo sedentario.

\textbf{¿Cómo lo hicieron?} \\
La revisión abarcó una variedad de estudios que utilizan sistemas wearables equipados con sensores de movimiento, como acelerómetros y giroscopios, que brindan retroalimentación en tiempo real sobre las posturas del cuerpo. Se evaluaron diferentes tipos de sistemas de retroalimentación, incluyendo vibraciones, alertas sonoras y visuales, y su efectividad para mejorar la postura en trabajadores de oficina y otros entornos laborales.

\textbf{Resultados:} \\
Se concluyó que los sistemas de retroalimentación biofísica son efectivos para reducir el riesgo de lesiones musculoesqueléticas, pero su efectividad varía según el tipo de retroalimentación utilizada y la duración del uso del dispositivo. Los resultados fueron más pronunciados en el corto plazo, aunque se señaló que los beneficios a largo plazo necesitan más investigación y validación. La retroalimentación táctil fue particularmente eficaz para mejorar la conciencia postural.

\subsection{Wearable Technology Assesses Surgeons' Posture During Surgery}
\textbf{Referencia:} Zulbaran-Rojas, Alejandro, et al. (2023) \cite{zulbaran2023wearable}\\

\textbf{¿Qué propusieron?} \\
Este estudio evaluó el uso de tecnología vestible para monitorear y evaluar la postura de los cirujanos durante procedimientos quirúrgicos. Dado el alto riesgo de trastornos musculoesqueléticos en cirujanos debido a las posturas estáticas prolongadas, este enfoque buscaba proporcionar datos objetivos para mejorar la ergonomía en el quirófano.

\textbf{¿Cómo lo hicieron?} \\
Equiparon a los cirujanos con sensores portátiles que registraban la postura y los movimientos durante las cirugías. Los sensores enviaban datos en tiempo real a un sistema de monitoreo que permitía la evaluación continua de la postura de los cirujanos. Este sistema fue probado en diferentes tipos de cirugía, incluyendo neurocirugía y cirugía general.

\textbf{Resultados:} \\
Los resultados mostraron que la tecnología vestible podía proporcionar una retroalimentación precisa sobre las posturas de los cirujanos, identificando posiciones de riesgo para la salud. Además, el monitoreo permitió a los cirujanos ajustar su postura durante los procedimientos, lo que podría ayudar a reducir el riesgo de lesiones musculoesqueléticas a largo plazo.

\subsection{Influence of the Wearable Posture Correction Sensor on Head and Neck Posture: Sitting and Standing Workstations}
\textbf{Referencia:} Ailneni, Ravi Charan, et al. (2019) \cite{ailneni2019influence}\\

\textbf{¿Qué propusieron?} \\
Este artículo evaluó cómo un sensor de corrección postural portátil influye en la postura de la cabeza y el cuello, tanto en estaciones de trabajo sentadas como de pie. El sensor proporciona retroalimentación en tiempo real cuando detecta una postura incorrecta, con el objetivo de mejorar la alineación de la cabeza y el cuello.

\textbf{¿Cómo lo hicieron?} \\
Los participantes usaron el sensor mientras realizaban tareas de trabajo en posiciones sentadas y de pie. Se registraron los ángulos de flexión y extensión del cuello, y el sistema proporcionó retroalimentación mediante una alerta vibratoria cuando se detectó una mala postura. Los participantes fueron evaluados antes y después del uso del sensor.

\textbf{Resultados:} \\
Los resultados mostraron una mejora significativa en la postura de la cabeza y el cuello cuando los participantes usaban el sensor, con una reducción del 20\% en los ángulos de flexión del cuello y un aumento en la conciencia postural. El sistema demostró ser útil para reducir el estrés en la columna cervical en entornos laborales.

\subsection{Effects of Using Wearable Devices on Reducing Sedentary Time and Prolonged Sitting in Healthy Adults: A Network Meta-Analysis}
\textbf{Referencia:} He, Zihao, et al. (2024) \cite{he2024effects}\\

\textbf{¿Qué propusieron?} \\
Este meta-análisis evaluó los efectos de los dispositivos portátiles en la reducción del tiempo sedentario y el tiempo de sentado prolongado en adultos sanos. Se analizaron diferentes tipos de intervenciones con dispositivos wearables, como relojes inteligentes y pulseras, para medir su efectividad en la promoción de la actividad física.

\textbf{¿Cómo lo hicieron?} \\
Se recopilaron y analizaron datos de múltiples estudios controlados aleatorios que evaluaban el uso de dispositivos portátiles para reducir el tiempo sedentario. El análisis incluyó datos de 26 estudios y 3,500 participantes, con una evaluación de los efectos de diferentes tipos de dispositivos y su duración de uso.

\textbf{Resultados:} \\
El meta-análisis reveló que los dispositivos portátiles fueron efectivos para reducir el tiempo sedentario y el tiempo de sentado prolongado, con una reducción promedio del 15-20\% en ambos. Los resultados sugieren que los wearables son herramientas útiles para promover la actividad física y mejorar los hábitos posturales, especialmente en poblaciones adultas.

\subsection{Real-Time Forward Head Posture Detection and Correction System Utilizing an Inertial Measurement Unit Sensor \cite{park2024real}}

\textbf{¿Qué hicieron?}\\

El estudio propuso un sistema en tiempo real para detectar y corregir la postura de cabeza adelantada (FHP, por sus siglas en inglés), un defecto postural comúnmente asociado con el dolor de cuello y una mala alineación espinal. El sistema utiliza un solo sensor de Unidad de Medición Inercial (IMU) que se coloca en el cuerpo del usuario para medir los cambios en la postura que podrían indicar el inicio de la FHP. El objetivo principal era proporcionar retroalimentación personalizada en tiempo real al usuario, alertándolo sobre la mala postura y guiándolo hacia una mejor alineación.

\textbf{¿Cómo lo hicieron?}\\

Los investigadores colocaron el sensor IMU en la columna cervical de los participantes y recolectaron datos en tiempo real sobre su postura del cuello. Usaron estos datos para determinar el Ángulo Cráneo-Vertebral (CVA), que es un indicador clave de la FHP. Se estableció un CVA base para cada participante y las desviaciones de esta base activaban alertas si el participante mantenía una mala postura durante más de 10 segundos. El sistema proporcionaba retroalimentación visual inmediata a través de una interfaz de usuario para fomentar la corrección de la postura. Se realizaron experimentos en un entorno controlado similar a una oficina con diez participantes masculinos, y se evaluó la capacidad del sistema para detectar y corregir la FHP.

\textbf{¿Cuáles fueron los resultados?}\\

Los resultados mostraron que el sistema fue efectivo para detectar la transición a la FHP, con un retraso de menos de 0.5 segundos. El sistema alertó con éxito a los participantes para que mantuvieran la postura correcta, lo que llevó a un aumento en su ángulo cráneo-vertebral (CVA) promedio. El sistema fue personalizado para cada participante, teniendo en cuenta sus condiciones físicas individuales. Cuando se proporcionó retroalimentación en tiempo real, los participantes mejoraron significativamente su postura, con un aumento promedio en su CVA de aproximadamente 3 a 5°. El sistema demostró un gran potencial para su aplicación en entornos donde los usuarios permanecen sedentarios por largos períodos, como el trabajo de oficina o el estudio.


\section{Metodología}

\subsection{Lista de componentes}
\subsubsection{Módulo MPU6050: Acelerómetro y giroscopio I2C}
\textbf{Función:} Medir la inclinación de la espalda mediante la detección de movimientos y rotaciones. \\
\textbf{Cantidad:} 1

\subsubsection{ESP32 DEVKIT V1 - NodeMCU 32 30-pin}
\textbf{Función:} Microcontrolador que se utilizará para procesar los datos del sensor y enviar la información de forma inalámbrica a la aplicación o servidor. \\
\textbf{Cantidad:} 2

\subsubsection{Pulsioxímetro MAX30102}
\textbf{Función:} Medir el nivel de oxigenación en sangre y la frecuencia del pulso. \\
\textbf{Cantidad:} 1

\subsubsection{Módulo Cargador de batería litio TP4056 con protección USB-C}
\textbf{Función:} Cargar la batería del dispositivo wearable. \\
\textbf{Cantidad:} 1

\subsubsection{Batería Plana 3.7V 450mAh 502040}
\textbf{Función:} Fuente de energía para los wearables. \\
\textbf{Cantidad:} 2

\subsubsection{Cable Jumper Dupont hembra a macho 10cm}
\textbf{Función:} Conectar los sensores y otros componentes al microcontrolador. \\
\textbf{Cantidad:} 1

\subsubsection{Módulo Buzzer Pasivo MH-FMD}
\textbf{Función:} Generar una alerta sonora en el dispositivo wearable para avisar al usuario sobre ciertos eventos. \\
\textbf{Cantidad:} 1

\subsection{Microcontrolador a usar}
El ESP32 será el microcontrolador encargado de procesar los datos de los sensores y comunicarse con la aplicación móvil a través de Bluetooth o Wi-Fi.

\subsection{Diagrama de bloques del sistema propuesto}
% Aquí incluiría un diagrama del sistema de sensores y microcontrolador

\section{Objetivos y Alcances}
\begin{itemize}
    \item \textbf{Objetivo General:} Desarrollar dos dispositivos wearables para medir la inclinación de la espalda y los niveles de oxigenación y pulso en tiempo real, utilizando sensores y tecnología IoT.
    
    \item \textbf{Objetivos Específicos:}
    \begin{itemize}
        \item Diseñar el circuito electrónico para medir la inclinación de la espalda utilizando el módulo MPU6050.
        \item Integrar el pulsioxímetro MAX30102 para medir el nivel de oxigenación en sangre y la frecuencia cardíaca del usuario.
        \item Desarrollar el firmware para el microcontrolador ESP32, encargado de procesar los datos de los sensores y enviarlos de manera inalámbrica.
        \item Construir una aplicación móvil que permita visualizar en tiempo real los datos obtenidos por los dispositivos wearables.
        \item Implementar una interfaz de usuario en la aplicación para la visualización de datos de inclinación, oxigenación y pulso.
        \item Garantizar la autonomía del dispositivo wearable mediante el uso de una batería recargable de litio, gestionada por un módulo TP4056.
        \item Incorporar una alerta sonora utilizando el buzzer para notificar al usuario sobre cambios importantes en los parámetros monitoreados.
    \end{itemize}

    \item \textbf{Alcances y limitaciones:}
    El alcance principal de este proyecto es proporcionar una solución en tiempo real para la medición y visualización de la inclinación de la espalda y los niveles de oxigenación y pulso. La visualización de los datos obtenidos de los dispositivos wearables se llevará a cabo únicamente cuando la aplicación esté abierta y en primer plano. En este sentido, la aplicación se encargará de mostrar los datos de manera activa mientras el usuario la esté utilizando. Sin embargo, en este primer ciclo de desarrollo, no se contempla la capacidad de la aplicación para realizar un seguimiento de los datos cuando se encuentre en segundo plano, lo que limita la funcionalidad en cuanto a la gestión y almacenamiento continuo de la información.

    Para una segunda iteración del proyecto, se plantea como objetivo la implementación de una funcionalidad adicional que permita a la aplicación almacenar y acceder a los datos históricos incluso cuando esté en segundo plano. Esto permitirá a los usuarios visualizar el historial de datos registrados en sesiones anteriores y realizar un análisis más detallado de su progreso a lo largo del tiempo.
\end{itemize} 

\bibliographystyle{ACM-Reference-Format}
\bibliography{bibitem}
\end{document}
\endinput
%%
%% End of file `sample-manuscript.tex'.
